\section{정리의 증명과 추가 실험}
\label{sec:proofs}

\subsection{Master Formula의 유도}
\label{sec:proof_master}

가우시안 소스 $k \sim \mathcal{N}(0, \Sigma)$에 대해, 블록-$b$ 직교 변환 후 독립 스칼라 $B$-bit 양자화의 왜곡을 유도한다.

\begin{proof}
\textbf{Step 1 (블록 대각 KLT).}
$\Sigma$를 $d/b$개의 $b \times b$ 블록으로 분할한다. 각 블록 $k$의 부분 공분산을 $\Sigma^{(k)}$라 하면, 블록 내 KLT는 각 블록의 고유벡터를 사용해 블록 내부를 decorrelate한다.

\textbf{Step 2 (고율 왜곡).}
고율 근사 아래에서, 분산이 $\sigma_j^2$인 가우시안 성분에 $b_j$-bit 스칼라 양자화를 적용하면 왜곡은 $D_j = c \cdot \sigma_j^2 \cdot 2^{-2b_j}$이다 \citep{cover2006elements}. 총 왜곡은 $D = \sum_j c \cdot \sigma_j^2 \cdot 2^{-2b_j}$이다.

\textbf{Step 3 (Water-Filling).}
$\sum_j b_j = d \cdot B$ 제약 하에서 라그랑지안 최적화를 적용하면, 최적 비트 배분은 $b_j^* = \alpha + \frac{1}{2}\log_2 \sigma_j^2$이고, 이때 총 왜곡은
\[
D^* = d \cdot c \cdot \left(\prod_{j=1}^{d} \sigma_j^2\right)^{1/d} \cdot 2^{-2B}.
\]

\textbf{Step 4 (블록 기하 평균).}
블록-$b$ 분할에서 각 블록의 기여는 $(\det \Sigma^{(k)})^{1/b}$이다. 차원 순열 $\pi$ 위에서 최적화하면
\[
G(b, \Sigma) = \min_\pi \left[\prod_{k=1}^{d/b} \det \Sigma^{(k)}_\pi\right]^{1/d}.
\]
따라서 $D^*(b, B, \Sigma) = d \cdot G(b, \Sigma) \cdot 2^{-2B}$.
\end{proof}

\subsection{블록 크기 계층 (따름정리~\ref{cor:block_hierarchy})의 증명}

\begin{proof}
$G(d) \leq G(1)$: 전체 $d$차원을 하나의 블록으로 보면 $G(d) = (\det \Sigma)^{1/d} = (\prod_i \lambda_i)^{1/d}$이고, 블록 크기 1에서는 $G(1) = (\prod_i \Sigma_{ii})^{1/d}$이다. Hadamard 부등식 $\det \Sigma \leq \prod_i \Sigma_{ii}$에 의해 $G(d) \leq G(1)$.

$G(1) \leq G_{\mathrm{uniform}}$: AM-GM 부등식에 의해 $(\prod_i \Sigma_{ii})^{1/d} \leq \frac{1}{d}\sum_i \Sigma_{ii} = \frac{\trace(\Sigma)}{d}$.

$G(d) \leq G(b)$: 블록-$d$ (전체 PCA)는 블록-$b$보다 더 큰 자유도를 가지므로 왜곡이 작거나 같다.
\end{proof}

\subsection{정리~\ref{thm:optimal_hamiltonian}의 증명}
\label{sec:proof_optimal}

\begin{proof}
위치 $\mathrm{pos}$에서 Post-RoPE 키 벡터 $k_{\mathrm{post}} = R_{\mathrm{RoPE}} k_{\mathrm{pre}}$의 공분산은
\[
\Sigma(\mathrm{pos}) = R_{\mathrm{RoPE}} \Sigma_0 R_{\mathrm{RoPE}}^\top
\]
이다.

\textbf{Step 1.} $\Sigma(\mathrm{pos})$의 고유분해는
\[
\Sigma(\mathrm{pos}) = (R_{\mathrm{RoPE}} V_0) \Lambda (R_{\mathrm{RoPE}} V_0)^\top
\]
이므로, 위치 $\mathrm{pos}$에서의 최적 회전은 $U^*(\mathrm{pos}) = R_{\mathrm{RoPE}}(\mathrm{pos}) \cdot V_0$이다.

\textbf{Step 2 (RoPE 상쇄).} Post-RoPE 키에 $U^*(\mathrm{pos})^\top$를 곱하면:
\[
U^*(\mathrm{pos})^\top k_{\mathrm{post}} = V_0^\top R_{\mathrm{RoPE}}^\top \cdot R_{\mathrm{RoPE}} k_{\mathrm{pre}} = V_0^\top k_{\mathrm{pre}}.
\]
\textbf{RoPE가 상쇄된다.} 따라서 위치에 무관하게 $V_0$를 Pre-RoPE 키에 적용하면 모든 위치에서 동시에 최적이 된다.

\textbf{Step 3 (이득 분석).} TurboQuant 대비 기대 MSE 이득의 상한은 Jensen 부등식에 의해 $R_{\mathrm{aniso}} = \mathrm{AM}(\lambda)/\mathrm{GM}(\lambda)$이다.
\end{proof}

\subsection{따름정리~\ref{cor:gain_over_turbo}의 증명}

\begin{proof}
Pre-RoPE PCA의 왜곡 계수는 $G_{\mathrm{PCA}} = \mathrm{GM}(\lambda)$이다. TurboQuant는 Haar 무작위 회전을 적용하며, 회전 후 $j$번째 성분의 기대 분산은 $\E[(R\Sigma_0 R^\top)_{jj}] = \mathrm{AM}(\lambda)$이다. Jensen 부등식에 의해
\[
\E[G_{\mathrm{TurboQuant}}] \leq \mathrm{AM}(\lambda).
\]
따라서 $\E[D_{\mathrm{TurboQuant}}] / D^*_{\mathrm{PCA}} \leq R_{\mathrm{aniso}}$.
\end{proof}

\subsection{Water-Filling floor=2의 실증적 동기}

$B$-bit 균일 양자화기의 MSE 비율 $c(B) = D(B) / \sigma^2$ (가우시안 소스, Max 1960):

\begin{center}
\begin{tabular}{@{}ccccc@{}}
\toprule
$B$ (bits) & 1 & 2 & 3 & 4 \\
\midrule
$c(B)$ & 0.3634 & 0.1175 & 0.03454 & 0.00942 \\
MSE 제거율 & 63.7\% & 88.3\% & 96.5\% & 99.1\% \\
\bottomrule
\end{tabular}
\end{center}

1-bit는 분산의 63.7\%만 제거하므로 비트 효율이 낮다. 그러나 이것은 \emph{MSE 관점}의 설명이다. Section~\ref{sec:exp_allocation_gap}의 E3b 실험에서 MSE 기준으로는 24/24 케이스 floor=0이 최적임이 확인되었으므로, floor=2의 진짜 이유는 MSE가 아닌 PPL에 있다. 이것은 Lloyd-Max 실패와 동일한 metric mismatch의 비트 배분 축 발현이다.

\subsection{5가지 가설 기각의 상세}
\label{sec:appendix_five_rejections}

\paragraph{가설 A 상세: Global $\kappa$ vs Per-head $\kappa$ spread.}
4모델에서 8개 레이어를 샘플링하여 per-head $\kappa$ 분포를 측정한 결과:

\begin{center}
\small
\begin{tabular}{@{}lcccc@{}}
\toprule
모델 & $\kappa$ median & $\kappa$ max & p95/median & Lloyd failure \\
\midrule
Qwen 1.5B & 64,127 & --- & --- & --- \\
Qwen 7B & 22,470 & --- & 5.0$\times$ & $1.05\times$ (mild) \\
Qwen 14B & 12,129 & --- & --- & --- \\
Mistral 7B & 14,321 & 2,805,296 & 15.3$\times$ & $5.06\times$ (severe) \\
\bottomrule
\end{tabular}
\end{center}

Global $\kappa$ median은 역상관(Qwen $>$ Mistral)이지만, per-head $\kappa$ spread (p95/median)는 $\rho = +1.0$으로 완벽 양상관한다. Mistral layer 2 H3의 $\kappa = 2,805,296$이 극단적 outlier이며, 이것이 per-layer breakdown (Section~\ref{sec:exp_per_layer})의 layer 2 최악 결과와 정확히 일치한다.

\paragraph{가설 B 상세: Hill estimator 측정.}
Top 10\% tail에서 Hill estimator로 측정한 tail index $\alpha$:

\begin{center}
\small
\begin{tabular}{@{}lccc@{}}
\toprule
모델 & $\alpha$ median & $\alpha$ min & 해석 \\
\midrule
Qwen 1.5B & 4.344 & 3.7 & Gaussian-like \\
Qwen 7B & 4.388 & 3.7 & Gaussian-like \\
Qwen 14B & 4.251 & 3.7 & Gaussian-like \\
Mistral 7B & 4.347 & 3.66 & Gaussian-like \\
\bottomrule
\end{tabular}
\end{center}

모든 모델 $\alpha \approx 4.3$으로, Gaussian-like 분포이다. Section~\ref{sec:theory_gap}의 $\kappa_4 \approx 0.5$ (mild non-Gaussianity) 결과와 일치한다.

\paragraph{가설 C 상세: Spherical quantizer.}
Mistral 64 heads에 Polar decomposition (3-bit angle + 1-bit magnitude) Spherical quantizer를 적용한 결과:

\begin{center}
\small
\begin{tabular}{@{}lccc@{}}
\toprule
Quantizer & MSE/Uniform (median) & Attn-MSE/Uniform & Win/64 \\
\midrule
$L^2$ Lloyd & 0.589 & --- & 64/64 \\
Spherical & 1.379 & 2.032 & 0/64 \\
\bottomrule
\end{tabular}
\end{center}

Spherical이 전면 실패한 이유: (1) k\_proj 출력에 RMSNorm 미적용, (2) Polar 분해가 key anisotropy를 포착 못 함, (3) 3-bit 각도 양자화가 coarse.

\paragraph{가설 D 상세: Fisher-Mahalanobis Lloyd.}
Per-token Fisher average로 whitening한 Mahalanobis Lloyd의 full-model PPL (Mistral-7B):

\begin{center}
\small
\begin{tabular}{@{}lcc@{}}
\toprule
Method & PPL & vs FP16 \\
\midrule
FP16 & 5.39\textsuperscript{$\dagger$} & --- \\
Uniform 2-bit & 243.38 & +4417\% \\
$L^2$ Lloyd 2-bit & 37.97 & +605\% \\
Fisher-Mahalanobis 2-bit & 356.40 & +6515\% \\
\bottomrule
\end{tabular}
\end{center}
\textsuperscript{$\dagger$}이 실험은 별도 calibration run (WikiText-2 2K 토큰 평가)에서 수행되어, 본문의 49K 토큰 평가와 수치가 다소 다르다 (FP16 5.39 vs 5.572). 방법 간 비율 비교는 유효하다.

Per-head Fisher-norm에서는 12/16 head에서 $L^2$ Lloyd를 이기지만, full-model PPL에서는 $9.4\times$ 악화된다. Information geometry(Fisher)가 PPL의 올바른 metric이 아님을 시사한다.

\paragraph{가설 E 상세: Discrete-WF.}
Gaussian single-channel에서 $r(b) = D_{\mathrm{uniform}}(b) / D_{\mathrm{Shannon}}(b)$:

\begin{center}
\small
\begin{tabular}{@{}lcccccc@{}}
\toprule
$b$ (bits) & 1 & 2 & 3 & 4 & 5 & 6 \\
\midrule
$r(b)$ & 1.46 & 1.91 & 2.41 & 2.98 & 3.62 & 4.36 \\
\bottomrule
\end{tabular}
\end{center}

$r(b)$는 단조 증가하며 $b=1$에서 knee가 없다. Heterogeneous WF 시뮬레이션(8 spectra $\times$ 3 budgets = 24 cases)에서 MSE-최적은 24/24 floor=0이다.

\subsection{Per-Layer Sensitivity 상세}
\label{sec:appendix_sensitivity}

\paragraph{Mistral-7B 전체 sensitivity ranking (32 layers).}

\begin{center}
\small
\begin{tabular}{@{}lr|lr|lr|lr@{}}
\toprule
Layer & $\Delta$PPL & Layer & $\Delta$PPL & Layer & $\Delta$PPL & Layer & $\Delta$PPL \\
\midrule
2 & +0.555 & 9 & +0.160 & 1 & +0.120 & 20 & +0.067 \\
4 & +0.521 & 22 & +0.155 & 30 & +0.116 & 21 & +0.067 \\
6 & +0.304 & 8 & +0.152 & 27 & +0.103 & 17 & +0.050 \\
3 & +0.287 & 23 & +0.122 & 29 & +0.096 & 15 & +0.047 \\
5 & +0.206 & 11 & +0.079 & 10 & +0.070 & 24 & +0.046 \\
7 & +0.166 &  &  &  &  &  & \\
\bottomrule
\end{tabular}
\end{center}

\paragraph{Qwen-7B 상위 sensitivity layers.}

\begin{center}
\small
\begin{tabular}{@{}lcccccc@{}}
\toprule
Layer & 0 & 26 & 4 & 5 & 22 & 나머지 (평균) \\
\midrule
$\Delta$PPL & +0.358 & +0.219 & +0.210 & +0.195 & +0.080 & $<$0.05 \\
\bottomrule
\end{tabular}
\end{center}

Mistral은 early layers (2--6)에 집중된 concentrated 패턴, Qwen은 early+late의 bimodal 패턴을 보인다. 이 차이가 sensitivity 배분의 model-specific 효과의 근본 원인이다.

\subsection{8가지 방법의 상세 해밀토니안 분류}
\label{sec:full_classification}

\begin{table}[h]
\centering
\caption{8가지 회전 기반 KV 캐시 양자화 방법의 해밀토니안 분류.}
\label{tab:hamiltonian_full}
\small
\begin{tabular}{@{}llcccl@{}}
\toprule
방법 & 학회 & 해밀토니안 $A$ & 블록 $b$ & 모드 & 비고 \\
\midrule
RoPE & --- & $A_{\mathrm{RoPE}}$ (고정) & 2 & I & 위치 인코딩 \\
TurboQuant \citep{zandieh2025turboquant} & ICLR'26 & $A \sim \mathrm{Haar}$ & $d$ & I & 등방화 \\
KVTC \citep{staniszewski2025kvtc} & ICLR'26 & $A_{\mathrm{PCA}}$ (공유) & $d$ & I & DP 비트 배분 \\
PolarQuant \citep{wu2025polarquant} & NeurIPS'25 & $A_{\mathrm{RoPE}}$ 구조 & 2 & II & 궤도 분해 \\
CommVQ \citep{li2025commvq} & ICML'25 & 중심화군 $C(T^{d/2})$ & 2 & I & RoPE 교환 \\
SpinQuant \citep{liu2024spinquant} & ICLR'25 & $A_{\mathrm{learn}}$ & $d$ & I & Cayley SGD \\
RotorQuant \citep{scrya2026rotorquant} & Preprint'26 & $\mathrm{Cl}(3,0)$ & 3 & I & Rotor \\
IsoQuant \citep{chen2026isoquant} & Preprint'26 & $\mathfrak{so}(4)$ & 4 & I & Isoclinic \\
\bottomrule
\end{tabular}
\end{table}

\subsection{추가 실험 상세}

\paragraph{MSE 순서 검증.}
3모델의 모든 KV 헤드(Qwen 112, Llama 256, Mistral 256, 총 624 헤드)에서 이론적 MSE 순서 $G(\mathrm{PCA}) \leq G(\mathrm{TurboQuant}) \leq G(\mathrm{uniform})$이 100\% 성립한다.

\paragraph{Lloyd-Max MSE 이득.}
Gaussian Lloyd-Max의 MSE 이득 (3모델 평균):

\begin{center}
\begin{tabular}{@{}lccc@{}}
\toprule
비트 & 2-bit & 3-bit & 4-bit \\
\midrule
Lloyd/Uniform MSE 비율 & 3.52$\times$ & 2.12$\times$ & 1.64$\times$ \\
\bottomrule
\end{tabular}
\end{center}

\paragraph{NIAH (Needle-In-A-Haystack).}
Qwen2.5-7B에서 4K/8K 컨텍스트 NIAH 검색 정확도:

\begin{center}
\begin{tabular}{@{}lcc@{}}
\toprule
방법 (2-bit) & 4K & 8K \\
\midrule
Pre-RoPE PCA & 100\% & 100\% \\
Random rot. & 100\% & 100\% \\
No rotation & 98\% & 94\% \\
\bottomrule
\end{tabular}
\end{center}
